\documentclass[10pt]{article}
\usepackage[letterpaper,margin=0.74in]{geometry}
\usepackage{amsmath,amssymb,booktabs,xcolor}
\usepackage[T1]{fontenc}
\usepackage[hidelinks]{hyperref}
\definecolor{ink}{HTML}{18364B}
\newcommand{\sect}[1]{\section*{\color{ink}#1}}
\newcommand{\code}[1]{\par\smallskip{\footnotesize\color{ink}\raggedright #1\par}}
\newcommand{\src}[1]{\texttt{\detokenize{#1}}}
\newcommand{\E}{\mathbb{E}}
\newcommand{\En}{\mathbb{E}_n}
\newcommand{\PP}{\mathbb{P}}
\newcommand{\cov}{\operatorname{Cov}}
\newcommand{\vech}{\operatorname{vech}}
\newcommand{\expit}{\operatorname{expit}}
\setlength{\parindent}{0pt}
\setlength{\parskip}{4pt}
\raggedbottom
\begin{document}
{\LARGE\bfseries DRFA simulation workflow}\par\smallskip
The original study ran 1,000 replications at each of \(n=250,500,1000\), with master seed \texttt{20261026}. The later replay used the same 1,000 saved \(n=500\) draws and compared two AIPW formulas under four nuisance conditions. The targets are the marginal covariance \(\Sigma=\cov(X)\) and 13 CFA parameters in \(\theta=(\Lambda,\Phi,\Omega)\).
\code{\src{python/drfa_simulation.py} lines 388--412, 457--469. Replay: \src{python/prepare_aipw_moment_comparison.py} lines 59--69, 90--102.}

\sect{1. Generate and mask the outcomes}
Draw \((X,Z)\) jointly Gaussian with all means 2, \(\operatorname{Var}(Z)=1\), and \(\cov(X_j,Z)=0.64\). Items 1--3 load on factor 1 and items 4--6 on factor 2. Each loading is \(0.8\), the factor variances are 1, their correlation is \(0.3\), and each uniqueness is \(0.36\). Thus \(\Sigma\) has diagonal 1, within-factor off-diagonal \(0.64\), and between-factor entries \(0.192\).
\[
\E[X_j\mid Z]=2+0.64(Z-2),\qquad
\cov(X_j,X_k\mid Z)=\Sigma_{jk}-0.64^2,\qquad
\pi_i=\PP(R_{ij}=1\mid Z_i)=\expit(Z_i-0.6).
\]
Draw \(R_{ij}\) independently across items given \(Z_i\). Mask \(X_{ij}\) when \(R_{ij}=0\). All four conditions use this same MAR mask and incomplete outcome matrix.
\code{\src{python/dgp.py} lines 6--7, 39--66. \src{python/drfa_simulation.py} lines 159--180.}

\sect{2. Fit the expectation and response models}
\begin{center}
\begin{tabular}{lcc}\toprule
Condition & Outcome design \(z_{E,i}\) & Response design \(z_{P,i}\)\\\midrule
Both correct & \((1,Z_i)\) & \((1,Z_i)\)\\
\(P\) misspecified & \((1,Z_i)\) & \((1,Z_i^2)\)\\
\(E\) misspecified & \((1,Z_i^2)\) & \((1,Z_i)\)\\
Both misspecified & \((1,Z_i^2)\) & \((1,Z_i^2)\)\\\bottomrule
\end{tabular}
\end{center}
Fit \(\widehat\beta_j\) by observed-item least squares and \(\widehat\alpha\) by Newton updates for the shared binomial logit using each row's observed-item count. Estimate \(\widehat\sigma_{jk}\) from the jointly observed fitted residuals, centering within each pair and dividing by its count minus one.
\begin{align*}
\widehat m_{ij}&=\widehat\E[X_{ij}\mid z_{E,i}]=z_{E,i}^{\top}\widehat\beta_j,
&\widehat\pi_i&=\widehat\PP(R_{ij}=1\mid z_{P,i})=\expit(z_{P,i}^{\top}\widehat\alpha),\\
\widehat q_{ijk}&=\widehat\E[X_{ij}X_{ik}\mid z_{E,i}]
=\widehat m_{ij}\widehat m_{ik}+\widehat\sigma_{jk}.
\end{align*}
\code{\src{python/drfa_simulation.py} lines 168--171. \src{python/aipw_m_estimation.py} lines 57--94, 259--315.\newline Lab notation: \src{aipw_m_estimation_and_cfa_variance_reference.tex} lines 20--49.}

\sect{3. Estimate the moments and form the CFA input}
Write \(e_{ij}=X_{ij}-\widehat m_{ij}\), \(w_{ij}=R_{ij}/\widehat\pi_i\), and let \(w_{ijk}=R_{ij}R_{ik}/\widehat\pi_i^2\) for \(j\ne k\) and \(w_{ij}\) for \(j=k\). Here \(\En\) averages over all rows.
\begin{align*}
\widehat\mu_j &= \En[\widehat m_{ij}+w_{ij}e_{ij}],\\
\widehat\nu_{jk}^{\rm raw} &= \En[\widehat q_{ijk}+w_{ijk}(X_{ij}X_{ik}-\widehat q_{ijk})],\\
\widehat\nu_{jk}^{\rm res} &= \En[\widehat q_{ijk}+\widehat m_{ij}w_{ik}e_{ik}
+\widehat m_{ik}w_{ij}e_{ij}+w_{ijk}(e_{ij}e_{ik}-\widehat\sigma_{jk})].
\end{align*}
The original formula weights the raw product correction jointly. The replay's residual version uses item weights for the linear residual terms and the joint weight for the residual product. Both pass \(S=\widehat Q=\widehat M-\widehat\mu\widehat\mu^\top\) to CFA, with \(\vech(\widehat M)=\widehat\nu\). They share \(\widehat\mu\) and the diagonal of \(S\). The single-misspecification arms probe double robustness, which requires correct \(P\) or both conditional moments on the \(E\) side.
\code{\src{python/aipw_m_estimation.py} lines 90--94, 329--337, 373--377. Residual version: \src{python/diagnose_aipw_covariance.py} lines 34--46. Pairing: \src{python/prepare_aipw_moment_comparison.py} lines 22--55.}

\newpage
\small
\setlength{\parskip}{3pt}
\sect{4. Population covariance and CFA input}
\[
\Sigma_X=\begin{bmatrix}
1 & .640 & .640 & .192 & .192 & .192\\
.640 & 1 & .640 & .192 & .192 & .192\\
.640 & .640 & 1 & .192 & .192 & .192\\
.192 & .192 & .192 & 1 & .640 & .640\\
.192 & .192 & .192 & .640 & 1 & .640\\
.192 & .192 & .192 & .640 & .640 & 1
\end{bmatrix},
\qquad
\begin{aligned}
S_{jk}
&=\widehat{\E}_{\rm AIPW}[X_jX_k]\\[-2pt]
&\quad-\widehat{\E}_{\rm AIPW}[X_j]\,\widehat{\E}_{\rm AIPW}[X_k].
\end{aligned}
\]
This fixed \(6\times6\) matrix is the target in every condition. For the joint Gaussian draw of \((X,Z)\), the DGP adds a seventh row and column with \(\cov(X_j,Z)=0.64\) and \(\operatorname{Var}(Z)=1\).
The optimizer receives \(S\) separately for each condition and replication. The row-level \(\widehat q_{ijk}\) contributes to the first term of this equation.
\code{\src{python/dgp.py} lines 39--53. \src{python/aipw_m_estimation.py} lines 329--337, 373--377. \src{python/torch_cfa_m_estimation.py} lines 76--80.}

\subsection*{\color{ink}5. Optimize the CFA model}
The covariance ML objective, up to constants and scaling, is
\[
L(\theta;S)=\tfrac12\bigl[\log|\Sigma_{\rm CFA}(\theta)|
+\operatorname{tr}\{S\Sigma_{\rm CFA}(\theta)^{-1}\}\bigr],
\qquad \Sigma_{\rm CFA}(\theta)=\Lambda\Phi\Lambda^\top+\Omega.
\]
Both engines fix cross-loadings to zero, set factor variances to one, and orient the first loading of each factor positive. The output is six loadings, six uniquenesses, and one factor correlation.

\textbf{Original PyTorch study.} Run L-BFGS with strong-Wolfe line search from three loading starts (\(0.8,0.6,1.0\)), initial uniquenesses \(0.4\), and correlation \(0\). Refine candidates by up to four Newton steps and require admissibility and agreement across converged starts. Use automatic differentiation of the CFA score to propagate the stacked AIPW sandwich to \(\widehat\theta\).
\code{\src{python/drfa.py} lines 216--221. \src{python/torch_cfa.py} lines 25--98, 110--139, 161--257. \src{python/torch_cfa_m_estimation.py} lines 76--128.}

\textbf{Paired lavaan replay.} Use the default NLMINB optimizer for covariance ML on the supplied \(S\), with no covariance rescaling and a 2,000-iteration limit. Skip non-positive-definite inputs, record nonconvergence and uniquenesses at or below \(10^{-4}\), and retain point estimates.
\code{\src{R/diagnose_aipw_covariance.R} lines 13--45. Optimizer default: lavaan 0.7-3, \src{lav_model_est}.}

\subsection*{\color{ink}Simulation statistics}
\begin{center}
\small
\begin{tabular}{lrrrr}\toprule
Replay condition & Invalid fits & Covariance MSE & CFA MSE & Common-valid draws\\
 & raw \(\to\) residual & raw \(\to\) residual & raw \(\to\) residual & \\\midrule
Both correct & \(20\to1\) & \(0.256\to0.161\) & \(0.121\to0.044\) & 979\\
\(P\) misspecified & \(12\to0\) & \(0.237\to0.141\) & \(0.112\to0.037\) & 988\\
\(E\) misspecified & \(112\to46\) & \(0.344\to0.306\) & \(0.148\to0.071\) & 863\\
Both misspecified & \(28\to0\) & \(0.243\to0.187\) & \(0.123\to0.047\) & 972\\\bottomrule
\end{tabular}
\end{center}
Covariance MSE uses all 1,000 draws and squared Frobenius loss. CFA MSE uses squared parameter error on the same admissible draws for both formulas. With both models misspecified, the covariance bias norm increased from \(0.103\) to \(0.238\) while empirical covariance variance decreased from \(0.233\) to \(0.130\). With \(E\) misspecified, non-positive-definite inputs increased from 27 to 29.
\code{\src{results/aipw_moment_comparison_n500}: \src{fit_summary.csv}, \src{statlab_validated_performance.csv}, \src{verification.json}. Metric code: \src{python/summarize_aipw_moment_comparison.py} lines 13--19, 33--69.}

\begin{center}
\small
\begin{tabular}{lccc}\toprule
Original study & \(n=250\) & \(n=500\) & \(n=1000\)\\
 & bias / coverage / invalid & bias / coverage / invalid & bias / coverage / invalid\\\midrule
Both correct & 1.75\% / 96.9\% / 196 & 1.26\% / 95.8\% / 20 & 0.58\% / 95.3\% / 2\\
\(P\) misspecified & 1.88\% / 96.9\% / 169 & 1.21\% / 95.9\% / 12 & 0.56\% / 95.2\% / 0\\
\(E\) misspecified & 2.26\% / 97.3\% / 370 & 1.53\% / 96.2\% / 114 & 0.88\% / 95.6\% / 23\\
Both misspecified & 4.19\% / 97.5\% / 252 & 3.76\% / 96.0\% / 29 & 3.27\% / 94.7\% / 0\\\bottomrule
\end{tabular}
\end{center}
Bias is relative \(L_2\) bias among admissible fits. Coverage averages parameter-level coverage of nominal 95\% sandwich Wald intervals among fits with usable variance estimates.
\code{\src{results/drfa_paper_study/n_*/summary.csv}. \src{python/drfa_simulation.py} lines 264--354.}

\end{document}
